\scanpage{203}
Il n'y a qu'un nb fini de pts de
$G_{\bar k}$ au dessus, \deleted{soit $U'$} la réunion des
de leurs composantes connexes est $U_{\bar k}$ [\uline{réf}]
\margin{$\times$}\deleted{donc $U$ \unsure{est q}} qui est de type fini sur
$\bar k$, donc $U$ l'est sur $k$. Comme $U_{\bar k}$ est
loc. irréductible, il en est de $=$ de $U$
\margin{$\times$}[réf]. Comme $e$ est rat/$k$, il s'ensuit
que sa composante connexe \inserted{\uline{\unsure{absolue}}} est définie sur $k$
[réf] donc \unsure{relt.|ablt.} irréductible; \deleted{et on voit} \inserted{Comme}
$G_{0\,\bar k}$ est invariant, $G_0$ est invariant.

\noindent\rule{1.15cm}{0.3pt}

Pour prouver l'existence du préschéma
quotient, on note plus généralement

\uline{Corollaire} \quad $G$ schéma en groupes loc de
type fini sur $k$, \deleted{\unsure{soit|G}}$H$ sous-schéma en
groupes ouvert \deleted{invariant} alors $H$ est
fermé, \deleted{Si $H$ est invariant, alors} $G/H$
existe, et $G\unsure{-|\to}G/H$ est fid. plat. (c'est
donc un \uuline{groupe} si $H$ est distingué)
% slanted
\margin{se généralise : \unsure{sur|une} base plus \unsure{général|générale}, quand $\deleted{\unsure{G}}H$ est \uuline{plat} sur \unsure{$k$|$S$} ? Oui si $S$ est \uuline{artinien}.}

\deleted{On peut utiliser une construction \unsure{spéciale} dans
le cas $H=G_0$.} Dans le cas $H=G_0$
on prend une composante connexe de $G$,
elle est stable sous $G_0$ [à signaler \ldots]
et \deleted{\unsure{puis|qui}} $G/G_0$ existe \deleted{\unsure{génériqut|géométriqut|groupe}}, car on le
voit trivialemt après l'extension finie de $k$.
% archive sheet number: 202
